% ==========================================================================
% Reference Mapping: Where Each Citation Appears in the Codebase
%
% This file tracks which references were used during implementation and
% WHERE they influenced design decisions. Use this as a lookup when
% writing the dissertation methodology and experimental setup chapters.
%
% Format: \section{Topic} -> bullet list of decisions with \cite{} refs.
% Compile with: pdflatex references.tex && bibtex references && pdflatex references.tex && pdflatex references.tex
% ==========================================================================

\documentclass[11pt,a4paper]{article}
\usepackage[utf8]{inputenc}
\usepackage[margin=2.5cm]{geometry}
\usepackage{natbib}
\usepackage{hyperref}
\usepackage{enumitem}

\title{Implementation Reference Map\\
\large Uncertainty-Conditioned RL with Evidential Policy Networks}
\author{Antonio Galdes}
\date{Last updated: \today}

\begin{document}
\maketitle

\section{Evidential Deep Learning Architecture}

\begin{itemize}[leftmargin=*]
  \item \textbf{NIG parameterisation} (gamma, nu, alpha, beta) with softplus constraints:
        \texttt{evidential\_policy.py}, \texttt{sb3\_integration.py}.
        Based on \cite{amini2020deep}.

  \item \textbf{Evidential regularisation} $L_R = |\mathbf{a} - \gamma| \cdot (2\nu + \alpha)$
        with $\lambda_\text{reg} = 0.01$:
        \texttt{sb3\_integration.py:EvidentialPPO.train()},
        \texttt{configs/train\_config.yaml}.
        Coefficient matches EPPO $\xi = 0.01$ \cite{eppo2024},
        validated range $[0.001, 0.1]$ by \cite{trustgeo2023} and \cite{soleimany2021evidential}.

  \item \textbf{Evidential actor (not critic)}: Design decision in \texttt{CLAUDE.md}.
        Note: EPPO \cite{eppo2024} applies evidential DL to the \emph{critic},
        not the actor. Our evidential actor is a less-studied configuration
        that produces Student-t policy distributions.

  \item \textbf{NIG hyperprior initialisation} (TODO in \texttt{train\_config.yaml}):
        nu bias $\approx 0.5$, alpha bias $\approx 0.5$, beta bias $\approx 0.0$.
        Produces ``honestly uncertain'' starting points.
        EPPO \cite{eppo2024} uses Gamma(2,1) for nu/alpha, Gamma(1,1) for beta.

  \item \textbf{Lambda annealing} from 0 to 0.01 over 50K steps (TODO):
        Lets NLL establish good predictions before regularisation begins.
        Standard practice in evidential DL literature.
\end{itemize}

\section{PPO Hyperparameters}

\begin{itemize}[leftmargin=*]
  \item \textbf{Learning rate 3e-4 with linear decay}:
        \texttt{train\_ppo.py:linear\_schedule()},
        \texttt{configs/train\_config.yaml}.
        3e-4 validated by \cite{andrychowicz2021what} (250K+ agents).
        Linear decay proven optimal by \cite{defazio2023road}.

  \item \textbf{batch\_size=256, n\_epochs=5} (40 off-policy steps, down from 320):
        \texttt{configs/train\_config.yaml}.
        \cite{carl2025} showed $\sim$15 off-policy steps outperformed 959 on CARLA.
        Corrected upward to 40 for single-env setup (no parallel data collection).

  \item \textbf{ent\_coef=0.0} (zero entropy for continuous control):
        \texttt{configs/train\_config.yaml}.
        \cite{andrychowicz2021what}: ``no evidence that the entropy term improves
        performance on continuous control environments.''

  \item \textbf{target\_kl=0.02} (early epoch stopping):
        \texttt{train\_ppo.py}, \texttt{configs/train\_config.yaml}.
        Standard PPO mechanism \cite{schulman2017proximal}. Set slightly higher
        than typical 0.015 due to noisy reward landscape from weather randomisation.

  \item \textbf{gamma=0.99} (kept, not reduced to 0.98):
        Scenario-specific: with current reward structure, success bonus at step
        $\sim$300 needs $\gamma^{300} \approx 0.05$ to propagate.
        At $\gamma = 0.98$, $\gamma^{300} \approx 0.002$ (invisible).
        SB3 highway-env parking uses $\gamma = 0.95$ \cite{highwayenv} but with
        dense reward per step.

  \item \textbf{n\_steps=2048} (kept for single CARLA env):
        $\sim$4 episodes per update provides trajectory diversity across weather
        conditions. \cite{carl2025} used lower values but with 4--8 parallel envs.

  \item \textbf{Observation normalisation} (VecNormalize, norm\_obs=True):
        \texttt{train\_ppo.py}.
        \cite{andrychowicz2021what}: ``one of the most impactful implementation details.''

  \item \textbf{norm\_reward=False}:
        \texttt{train\_ppo.py}.
        Auto-normalisation divides by $\text{std}(\text{discounted returns})$,
        unstable when reward distribution shifts (early: penalties only, later:
        success bonuses). With potential-based shaping, manual scaling preferred.
\end{itemize}

\section{Network Architecture}

\begin{itemize}[leftmargin=*]
  \item \textbf{net\_arch=[256,256]} for both actor and critic:
        \texttt{configs/train\_config.yaml}.
        Evidential head outputs $4 \times 3 = 12$ NIG parameters from 18-dim state.
        \cite{andrychowicz2021what} recommend separate sizing; kept equal because
        the 18-dim state includes 9 covariance features encoding complex non-linear
        uncertainty structure, plus 3 relative target pose features.

  \item \textbf{ReLU activation}:
        Both EPPO \cite{eppo2024} and \cite{amini2020deep} use ReLU.

  \item \textbf{LayerNorm} (TODO):
        EPPO \cite{eppo2024} uses LayerNorm after each hidden layer.
        Preserves network plasticity under non-stationary training
        \cite{dohare2024loss, lyle2022understanding}.
\end{itemize}

\section{Reward Function}

\begin{itemize}[leftmargin=*]
  \item \textbf{Potential-based distance shaping} (TODO, Task 9):
        Switch from $R = -d$ to $R = d_{t-1} - d_t$.
        Grounded in \cite{ng1999policy} (policy invariance under
        potential-based transformations).

  \item \textbf{Milestone bonuses} at distance thresholds (TODO):
        \cite{rarlap2024} found milestone-augmented rewards outperformed
        both sparse and pure dense rewards for parking with PPO.

  \item \textbf{highway-env parking benchmark} as reference design:
        \cite{highwayenv} uses weighted p-norm with $p = 0.5$,
        weights $[1.0, 0.3, 0.0, 0.0, 0.02, 0.02]$.
\end{itemize}

\section{Environment Design / Curriculum and Bay Sampling}

\begin{itemize}[leftmargin=*]
  \item \textbf{Stratified bay type sampling} (1/3 each: perpendicular, angled, parallel):
        \texttt{carla\_parking.py:CARLAParkingEnv.reset()},
        \texttt{configs/train\_config.yaml:bay\_type\_sampling}.
        Justified on three grounds:
        \begin{itemize}
          \item \cite{bengio2009curriculum}: balanced difficulty distribution leads to
                faster convergence than training on the natural (skewed) distribution.
                The parking lot has far more perpendicular bays than parallel bays
                in the wild; without stratification the agent would rarely see the
                hardest manoeuvre.
          \item \cite{florensa2017reverse}: rare difficult scenarios (parallel pull-in bays)
                are never practised enough without deliberate curriculum design. The
                stratified sampler is the minimal curriculum intervention that guarantees
                each bay type is seen equally often.
          \item \cite{leurent2018highway} highway-env ParkingEnv uses 28 same-type slots
                as the direct precedent for a goal-conditioned parking environment.
                This environment extends the design to 15 bays across 3 types with
                stratified sampling.
          \item \cite{cobbe2019quantifying}: general instance diversity principle for RL
                generalisation. Cited only for the general principle -- the specific
                count thresholds (4000--16000 CoinRun levels) do not transfer to
                vector-state parking, where the image-based pixel-memorisation failure
                mode does not apply.
        \end{itemize}
\end{itemize}

\section{Environment Design and Domain Randomisation}

\begin{itemize}[leftmargin=*]
  \item \textbf{Randomised floor plan presets (2 training, 1 OOD holdout)}:
        \texttt{configs/train\_config.yaml:parking\_scenarios.floor\_plans}.
        Two floor plan geometries (\texttt{rectangle}, \texttt{trapezoid}) are
        sampled uniformly during training. One geometry (\texttt{irregular\_a},
        five-sided polygon) is held back entirely and used exclusively for OOD
        evaluation.
        Justified by domain randomisation \cite{tobin2017domain}: exposing the
        policy to geometric variety during training encourages generalisation.
        The held-out geometry drives epistemic uncertainty higher on novel
        layouts, demonstrating the evidential policy's calibration.
        Why 3 floor plans rather than more: the agent's relative-obs formulation
        \texttt{(dx, dy, dyaw)} makes absolute bay position in world coordinates
        invisible -- overfitting to specific bay positions is structurally
        impossible. The real diversity axes are manoeuvre type (3 bay types) and
        lot perimeter shape. Two training floor plans provide sufficient geometric
        variety. See also \cite{cobbe2019quantifying} for the general diversity
        principle (note: their 4000--16000 CoinRun level counts do not transfer
        to vector-state parking -- cited for the general principle only).

  \item \textbf{Randomised bay occupancy per episode} ($\sim$70\% probability
        per eligible slot): Prevents the policy from memorising a fixed obstacle
        layout. The target slot and parallel-bay neighbours are always kept empty.
        Also justified by domain randomisation \cite{tobin2017domain}.

  \item \textbf{Randomised GNSS noise tier, NPC patrol count, pedestrian count per episode}:
        Varies the localisation uncertainty distribution across episodes.
        Per-episode RTK fix-state tier sampled from
        \texttt{configs/gnss\_noise\_profiles.yaml}. 0--3 patrol vehicles
        (scripted proportional controller, no Traffic Manager), 0--4
        pedestrians (\texttt{WalkerControl}, zone-constrained random-walk,
        no NavMesh required). No weather variation (FlatPlane does not
        render weather effects). \cite{tobin2017domain}.

  \item \textbf{Pedestrian zone-boundary correction (steer toward zone centre)}:
        \texttt{carla\_parking.py:\_update\_pedestrians()}.
        Pedestrians are assigned to aisle zones defined in the layout YAML.
        When a walker reaches within 0.5\,m of the zone boundary, it is assigned
        a new heading directed toward the zone centre rather than using velocity
        reflection.
        Pure reflection causes oscillation: the walker overshoots by one physics
        step, exits the zone, the heading is negated, and on the next step the
        walker is still outside and is negated again, producing a standing
        vibration on the boundary wall.
        The ``steer toward interior'' approach is the same mechanism used by the
        CARLA Scenario Runner \cite{carla_scenario_runner}, which assigns a new
        interior waypoint when an actor approaches a trigger-region boundary.
        Pedestrians remain in the aisles between parked car bodies, which is
        the correct location for LiDAR occlusion effects.

  \item \textbf{Three bay types (perpendicular, angled 45 deg, parallel forward entry)}:
        \texttt{configs/layouts/*.yaml:bays[*].bay\_type}.
        All three bay types appear in real car parks, including within the same
        facility \cite{ced_parking_design}. Training across all three bay types
        improves generalisation \cite{seg_parking2024}. Bay type is not in the
        observation; the agent infers the required manoeuvre from the relative
        target pose \texttt{(dx, dy, dyaw)} and reward signal.

  \item \textbf{Bay dimensions from German EAR 05 / EU harmonised practice (FGSV 2005)}:
        Perpendicular 2.5 m $\times$ 5.0 m, aisle 6.0 m;
        angled 2.5 m $\times$ 5.4 m, aisle 3.6 m;
        parallel 2.5 m $\times$ 8.0 m (forward pull-in without reversing), aisle 4.0 m.
        EAR 05 represents current comfortable minimum standards, slightly larger
        than absolute minimums, appropriate for a learning agent.
        Source: FGSV Empfehlungen fur Anlagen des ruhenden Verkehrs (EAR 05), 2005.

  \item \textbf{Reserved slots for parallel forward entry}:
        Slots immediately in front and behind the target parallel bay are marked
        \texttt{always\_empty} in the layout YAML
        (\texttt{configs/layouts/*.yaml:bays[*].always\_empty: true}).
        Engineering decision: no reference required. Perpendicular and angled
        bays require no reserved neighbours as the ego drives straight in without
        sweeping adjacent slots.
\end{itemize}

\section{Experimental Design and Evaluation}

\begin{itemize}[leftmargin=*]
  \item \textbf{10 seeds minimum} for final ablation:
        \texttt{configs/train\_config.yaml:random\_seeds}.
        \cite{henderson2018deep} showed 5 seeds produce misleading results.
        \cite{colas2018many} showed 5 seeds yields insufficient statistical power.
        \cite{agarwal2021deep} established IQM + bootstrap CIs as gold standard.

  \item \textbf{rliable evaluation protocol}:
        Report IQM with 95\% stratified bootstrap CIs, probability of improvement,
        performance profiles. \cite{agarwal2021deep}.

  \item \textbf{Sequential seeds [0..9]}:
        Avoids appearance of cherry-picking. Standard practice per
        \cite{agarwal2021deep}.

  \item \textbf{1M total timesteps}:
        EPPO \cite{eppo2024} uses 500K. \cite{lazzaroni2024parking} used 60M
        with curriculum learning. 1M provides 2$\times$ safety margin for
        compact state + shaped rewards.
\end{itemize}

\section{System Framing and Sensor Architecture}

\begin{itemize}[leftmargin=*]
  \item \textbf{Decoupled perception/control architecture}:
        Target parking pose $(x, y, \theta)$ is provided by an upstream visual
        perception module (e.g.\ camera-based bay detector). This system focuses
        on the control layer only: navigating safely to the given pose under
        localisation uncertainty. This matches commercial autonomous parking
        pipelines \cite{bosch_automated_parking, valeo_park4u} where perception
        and control are decoupled. Onboard bay detection is out of scope and
        noted as future work in the dissertation.

  \item \textbf{RTK-GNSS for localisation (outdoor lots only)}:
        RTK-GNSS provides cm-level accuracy under nominal conditions (RTK fixed),
        with well-understood degradation modes (float $\sim$30\,cm, standalone
        $\sim$2--5\,m) that produce the covariance variation the policy needs.
        A base station + rover costs $\sim$600\,GBP and is trivially deployable
        at an outdoor lot. Parking bays are 2.5\,m wide with $\sim$0.35\,m margin
        per side; sub-0.3\,m lateral accuracy is required \cite{reid2019localization}.
        RTK fixed satisfies this; standalone does not.
        Justified by \cite{takasu2009rtklib} (cm-level with low-cost receivers),
        \cite{groves2013gnss} (multipath/NLOS degradation in structured environments).
        Scope limitation: outdoor parking lots only (stated in dissertation).

  \item \textbf{Why RTK-GNSS replaced Cartographer scan-matching}:
        The earlier approach used 2D LiDAR scan-matching via Cartographer
        \cite{cartographer2016}. This was abandoned because: (1) when the lot
        is full, parked cars occlude perimeter cones and localisation degrades
        catastrophically; (2) the Cartographer stack added massive complexity
        (pbstream mapping workflow, Lua configs, TfToOdomNode, IMU frame relay,
        cross-distro DDS hacks). The thesis contribution is uncertainty-conditioned
        RL, not the localisation method.

  \item \textbf{Wheel odometry excluded from EKF sensor suite}:
        The CARLA ROS bridge publishes \texttt{/carla/ego\_vehicle/odometry}
        from ground-truth physics state (not integrated wheel encoder ticks).
        It carries zero noise and zero covariance, so feeding it into the EKF
        suppresses covariance variation regardless of lot conditions.
        Prediction step uses IMU only between GNSS fixes.
        Justified by \cite{moore2014generalized}: the \texttt{robot\_localization}
        package supports IMU-only prediction; wheel odometry is optional.
        Parking manoeuvres are short ($\sim$10--30\,s) and low-speed
        ($<$2\,m/s). IMU-only drift over this duration is bounded and
        corrected at each GNSS fix.

  \item \textbf{Sensor architecture (single suite)}:
        \begin{itemize}
          \item RTK-GNSS (base + rover): position fixes for EKF localisation
                ($\sim$600\,GBP). \cite{takasu2009rtklib, reid2019localization}.
          \item IMU: orientation, yaw rate for EKF prediction step ($\sim$50\,GBP).
          \item 2D LiDAR: obstacle clearance only (obs indices 15--19), not used
                for localisation ($\sim$500--1500\,GBP).
                See \cite{tang2015lidar, niu2017online, chen2020lidar} for
                IMU+LiDAR sufficiency in parking scenarios.
        \end{itemize}

  \item \textbf{RTK fix states as natural uncertainty source}:
        No prior work conditions RL vehicle control on GNSS fix-state covariance.
        The RTK fix-state tiers (fixed $\sim$2\,cm, float $\sim$30\,cm,
        standalone $\sim$2--5\,m) create natural discrete uncertainty regimes
        that the policy must learn to handle. These tiers are modelled in
        \texttt{configs/gnss\_noise\_profiles.yaml} and sampled per-episode
        during training. \cite{reid2019localization, groves2013gnss}.

  \item \textbf{Functional safety justification}:
        ISO 26262 \cite{iso26262} requires systems to detect and respond to
        faults, not just operate correctly under nominal conditions. The
        uncertainty-aware policy satisfies this: it detects localisation
        degradation (via EKF covariance) and responds safely (via evidential
        epistemic uncertainty triggering handoff).

  \item \textbf{ODD boundary detection via epistemic uncertainty}:
        SAE J3016 \cite{sae_j3016} defines ODD: conditions under which an
        automated system is designed to function. Epistemic uncertainty from
        the evidential actor serves as an ODD boundary detector -- high
        epistemic indicates operation outside the trained envelope.
\end{itemize}

\section{Infrastructure}

\begin{itemize}[leftmargin=*]
  \item \textbf{CARLA simulator}: \cite{dosovitskiy2017carla}.
        Headless, vector observations only (no camera rendering).

  \item \textbf{Stable-Baselines3}: \cite{raffin2021stable}.
        PPO implementation, VecNormalize, callbacks, custom policy subclassing.

  \item \textbf{robot\_localisation EKF}: \cite{robotlocalization, moore2014generalized}.
        Fuses RTK-GNSS position fixes (via \texttt{navsat\_transform\_node})
        with IMU orientation/yaw rate. Provides covariance output as
        uncertainty signal in the observation space.

  \item \textbf{Cartographer (historical)}: \cite{cartographer2016}.
        HISTORICAL: was used for LiDAR scan-matching in an earlier project
        iteration. Replaced by RTK-GNSS+IMU because cone occlusion by parked
        cars caused catastrophic scan-match failure in full lots.
\end{itemize}

\section{Novelty Positioning}

\begin{itemize}[leftmargin=*]
  \item \textbf{Evidential actor (vs EPPO)}: EPPO \cite{akgul2025eppo} applies
        EDL to the critic V(s) only; actor remains standard Gaussian. We place
        the NIG evidential head on the actor network itself, yielding per-action
        epistemic and aleatoric uncertainty in a single forward pass.

  \item \textbf{EDL on Q-network (discrete)}: CEQR-DQN \cite{stutts2024ceqr}
        applies EDL to the Q-network for MinAtar (discrete actions). Different
        problem class from continuous-action parking control.

  \item \textbf{EDL on perception (not policy)}: DRO-EDL-MPC \cite{ham2025dro}
        applies EDL to perception (object detection) to inform MPC constraints.
        Not RL. We place EDL at the policy level within model-free RL.

  \item \textbf{Ensemble uncertainty conditioning}: UNICON \cite{yu2022unicon}
        uses ensemble Q-disagreement to condition policy selection. Not evidential
        (ensemble-based). MuJoCo benchmarks, not driving/parking.

  \item \textbf{Pedestrian prediction covariance in state}: Cai et al.\
        \cite{cai2024crowd} feed pedestrian trajectory prediction covariance
        into DRL state. Different uncertainty source (other agents, not ego
        localisation) and domain (crowd navigation, not parking).

  \item \textbf{Parking uncertainty gap confirmed}: A 2025 MDPI Sensors survey
        \cite{parking_survey2025} identifies uncertainty estimation as critical
        for parking safety but notes it exists only at perception level, not
        integrated into control/decision-making policy.
\end{itemize}

\bibliographystyle{plainnat}
\bibliography{references}

\end{document}
